\chapter{QED for Hadrons}
\label{ch:qed-for-hadrons}

\section{Motivation: QED probes internal structure}
In Quantum Electrodynamics (QED), elementary charged fields couple to the photon
through the conserved electromagnetic (e.m.) current. For a set of elementary
Dirac fermions $\psi_j$ with charges $q_j$,
\begin{equation}
\mathcal L_{\rm QED}
= -\frac14 F_{\mu\nu}F^{\mu\nu}
+ \sum_{j=1}^n \bar\psi_j \bigl(i\gamma^\mu D_\mu^{(j)} - m_j\bigr)\psi_j,
\qquad
D_\mu^{(j)}=\partial_\mu+i q_j A_\mu .
\end{equation}
The interaction term is
\begin{equation}
\mathcal L_I = -\sum_{j=1}^n q_j\,\bar\psi_j\gamma^\mu A_\mu \psi_j
\equiv -\sum_{j=1}^n j_j^\mu A_\mu \equiv - j^\mu A_\mu,
\qquad \partial_\mu j^\mu=0.
\end{equation}
If the external states are elementary, amplitudes are computable perturbatively in $\alpha$.

\medskip
\noindent\textbf{Hadrons are not elementary:} they are composite bound states of quarks and gluons.
At sufficiently low momentum transfer the photon sees the hadron as point-like,
but as the momentum transfer increases, the photon resolves the hadron's internal
structure. We typically cannot compute this structure from first principles
without doing QCD. Instead, we \textbf{parameterize our ignorance} in the most general way allowed by:
\begin{itemize}
  \item space-time translation invariance,
  \item Lorentz covariance,
  \item discrete symmetries ($P$, $C$ when appropriate),
  \item and current conservation (gauge invariance).
\end{itemize}
The result is the notion of \textbf{form factors} (for exclusive processes) and
\textbf{structure functions} (for inclusive processes).

\section{Form factors for composite hadrons}
\subsection{General logic: what you can write down}
Suppose an electromagnetic process involves a hadronic matrix element of the QCD
electromagnetic current
\begin{equation}
j^\mu_{\rm QCD}(x)=\sum_{f=u,d,s,\dots} q_f\,\bar q_f(x)\gamma^\mu q_f(x).
\end{equation}
The object we do \emph{not} know how to compute in general is a matrix element such as
$\langle {\rm hadrons} | j^\mu_{\rm QCD}(x) | {\rm hadrons} \rangle$.
We proceed in steps:

\begin{enumerate}
  \item \textbf{Translation invariance} fixes the $x$-dependence:
  \[
  \langle f|j^\mu_{\rm QCD}(x)|i\rangle
  = e^{-i(p_f-p_i)\cdot x}\,\langle f|j^\mu_{\rm QCD}(0)|i\rangle .
  \]
  \item \textbf{Lorentz covariance} constrains the tensor/spinor structures that can appear.
  \item \textbf{Discrete symmetries} remove forbidden structures (e.g.\ parity eliminates certain $\gamma_5$ terms in vector-current matrix elements between positive-parity states).
  \item \textbf{Current conservation} ($\partial_\mu j^\mu=0$) implies $q_\mu \langle f|j^\mu|i\rangle=0$ and further reduces the number of independent terms.
\end{enumerate}

All remaining freedom is encoded in a \emph{finite number} of scalar functions of
Lorentz invariants (e.g.\ $q^2$), called \textbf{form factors}.

\subsection{Pion electromagnetic form factor}
Consider $e^-e^+\to \pi^+\pi^-$. In the center-of-mass frame, at energies where the pion
structure starts to matter, the pion couples through the quark current $j^\mu_{\rm QCD}$.

The unknown hadronic ingredient is
\begin{equation}
\langle \pi^+(p_2)\,\pi^-(p_1)|\,j^\nu_{\rm QCD}(x)\,|0\rangle.
\end{equation}
Translation invariance gives
\begin{equation}
\langle \pi^+\pi^-|\,j^\nu_{\rm QCD}(x)\,|0\rangle
= e^{-i(p_1+p_2)\cdot x}\,
\langle \pi^+\pi^-|\,j^\nu_{\rm QCD}(0)\,|0\rangle .
\end{equation}
Lorentz covariance implies the most general decomposition (with scalar functions of $s=(p_1+p_2)^2$):
\begin{equation}
\langle \pi^+(p_2)\,\pi^-(p_1)|\,j^\nu_{\rm QCD}(0)\,|0\rangle
= e\,(p_1+p_2)^\nu \,\widetilde F_\pi(s)\;+\;e\,(p_2-p_1)^\nu\,F_\pi(s).
\end{equation}
Now impose current conservation:
\begin{equation}
0=(p_1+p_2)_\nu \langle \pi^+\pi^-|\,j^\nu_{\rm QCD}(0)\,|0\rangle
= e\,(p_1+p_2)^2\,2\,\widetilde F_\pi(s),
\end{equation}
hence $\widetilde F_\pi(s)=0$. Therefore,
\begin{equation}
\boxed{
\langle \pi^+(p_2)\,\pi^-(p_1)|\,j^\nu_{\rm QCD}(0)\,|0\rangle
= e\,(p_2-p_1)^\nu\,F_\pi(s)
}
\end{equation}
where $F_\pi(s)$ is the \textbf{pion electromagnetic form factor}. At low energies
(near threshold) the pion behaves approximately point-like, so
\begin{equation}
F_\pi(s)\simeq 1 \qquad \text{for } s \gtrsim 4m_\pi^2 \text{ and small momentum transfer}.
\end{equation}
Experiment measures $F_\pi(s)$; in principle QCD can predict it.

\subsection{Spin-$\tfrac12$ hadrons: proton form factors and vertex decomposition}
For elastic electron--proton scattering $e^-p\to e^-p$, the hadronic ingredient is
\begin{equation}
\langle p(p_2,\lambda_2)|\,j^\nu_{\rm QCD}(x)\,|p(p_B,\lambda_B)\rangle .
\end{equation}
Translation invariance yields
\begin{equation}
\langle p(p_2)|\,j^\nu_{\rm QCD}(x)\,|p(p_B)\rangle
= e^{-i(p_2-p_B)\cdot x}\,
\langle p(p_2)|\,j^\nu_{\rm QCD}(0)\,|p(p_B)\rangle ,
\end{equation}
so define the momentum transfer $q=p_2-p_B$, $q^2=t$.

Lorentz covariance allows the most general Dirac decomposition consistent with a vector current.
After imposing parity invariance (eliminating $\gamma_5$ structures) and current conservation
($q_\nu \langle p|j^\nu|p\rangle=0$), the matrix element can be written in terms of two
independent scalar functions:
\begin{equation}
\boxed{
\langle p(p_2,\lambda_2)|\,j^\nu_{\rm QCD}(0)\,|p(p_B,\lambda_B)\rangle
= \bar u_{\lambda_2}(p_2)\left[
e\,F_1(q^2)\gamma^\nu
+ e\,\frac{F_2(q^2)}{2m_p}\,i\sigma^{\nu\rho}q_\rho
\right]u_{\lambda_B}(p_B).
}
\label{eq:dirac-pauli-ff}
\end{equation}
Here:
\begin{itemize}
  \item $F_1(q^2)$ is the \textbf{Dirac form factor} (charge distribution),
  \item $F_2(q^2)$ is the \textbf{Pauli form factor} (magnetization/anomalous moment),
  \item $\sigma^{\nu\rho}=\frac{i}{2}[\gamma^\nu,\gamma^\rho]$.
\end{itemize}

\paragraph{Low-$q^2$ normalization and anomalous magnetic moment.}
If the proton were a minimally coupled point Dirac particle, one would expect
$F_1(0)=1$ and $F_2(0)=0$. Experiment instead finds a large anomalous magnetic moment.
In an effective-field-theory spirit, one can include the leading non-minimal (dimension-5)
Pauli term compatible with gauge invariance, Lorentz invariance and discrete symmetries:
\begin{equation}
\delta\mathcal L
= \frac{e\,\kappa_p}{2m_p}\,\bar\psi_p \sigma^{\mu\nu}F_{\mu\nu}\psi_p,
\qquad \kappa_p=\frac{g_p}{2}-1 \simeq 1.79,\quad g_p\simeq 5.58.
\end{equation}
Then the low-$q^2$ behavior becomes
\begin{equation}
F_1(0)=1,\qquad F_2(0)=\kappa_p.
\end{equation}
Similarly for the neutron: $F_1(0)=0$ (net charge $0$) and $F_2(0)=\kappa_n\simeq -1.91$.

\subsection{Unpolarized elastic scattering: hadronic tensor and Sachs form factors}
If beams and detectors are unpolarized, we do not need the full spin-dependent matrix element,
only the spin-averaged product (hadronic tensor):
\begin{equation}
L^{\mu\nu}_p(p_B,p_2)
\equiv \frac12 \sum_{\lambda_B,\lambda_2}
\langle p(p_2,\lambda_2)|j^\mu_{\rm QCD}(0)|p(p_B,\lambda_B)\rangle\,
\langle p(p_2,\lambda_2)|j^\nu_{\rm QCD}(0)|p(p_B,\lambda_B)\rangle^\ast .
\end{equation}
Lorentz and parity invariance imply that $L^{\mu\nu}_p$ can be decomposed into a small
basis of tensors built from $g^{\mu\nu}$, $p^\mu=(p_B+p_2)^\mu$, and $q^\mu=(p_2-p_B)^\mu$.
Current conservation removes terms proportional to $q^\mu$ in the appropriate way, leaving
\begin{equation}
L^{\mu\nu}_p
= 2m_p^2\,G_1(q^2)\left(-g^{\mu\nu}+\frac{q^\mu q^\nu}{q^2}\right)
+ G_2(q^2)\,p^\mu p^\nu,
\end{equation}
so unpolarized elastic scattering depends on \textbf{two} scalar functions, consistent
with the two form factors in \eqref{eq:dirac-pauli-ff}.

It is conventional (and physically transparent) to use the \textbf{Sachs form factors}
$G_E$ and $G_M$, defined (with $Q^2\equiv -q^2>0$ and $\tau\equiv Q^2/(4m_p^2)$) by
\begin{equation}
G_E(Q^2)=F_1(Q^2)-\tau F_2(Q^2),
\qquad
G_M(Q^2)=F_1(Q^2)+F_2(Q^2).
\end{equation}
Interpretation:
\begin{itemize}
  \item $G_E$ is associated with the \textbf{electric charge distribution},
  \item $G_M$ is associated with the \textbf{magnetization distribution}.
\end{itemize}
For a point-like minimally coupled Dirac particle one expects $G_E=G_M=1$.

\paragraph{Rosenbluth-type form in the LAB frame.}
Early electron--proton scattering experiments were performed in fixed-target kinematics (LAB).
The differential cross section can be organized so that $G_E$ and $G_M$ may be extracted
from angular dependence (Rosenbluth separation). In the notes' conventions one often finds a form
\begin{equation}
\left(\frac{d\sigma}{d\Omega}\right)_{\rm LAB}
= \frac{\alpha^2 E_1}{4E_A^3\sin^4(\theta/2)}
\left(G_2\cos^2\frac{\theta}{2}+G_1\sin^2\frac{\theta}{2}\right),
\end{equation}
with $G_1,G_2$ re-expressed in terms of $G_E,G_M$ and kinematic factors.
Empirically, a simple approximate fit is
\begin{equation}
G_E(Q^2)\simeq \frac{1}{1+Q^2/Q_0^2},
\qquad
G_M(Q^2)\simeq \frac{g_p}{2}\,G_E(Q^2),
\qquad
Q_0^2\simeq 0.71~{\rm GeV}^2,
\end{equation}
illustrating that form factors decrease with momentum transfer: larger $Q^2$ resolves
smaller distance scales and reveals finite size.

\section{Deep Inelastic Scattering (DIS): from form factors to structure functions}
\subsection{Kinematics and why DIS is different}
Consider inclusive scattering
\[
e^-(p_A)+p(p_B)\to e^-(p_1)+X,
\]
where $X$ is \emph{any} hadronic final state. We measure only the outgoing electron.
Define the momentum transfer $q=p_A-p_1$ and $Q^2=-q^2>0$.

Compared with elastic scattering (where $X=p$), in DIS we must sum over all hadronic states:
\begin{equation}
\sum_X (2\pi)^4 \delta^{(4)}(p_A+p_B-p_1-p_X)\,
\langle X|j^\mu_{\rm QCD}(0)|p\rangle
\langle X|j^\nu_{\rm QCD}(0)|p\rangle^\ast.
\end{equation}
This motivates defining the \textbf{hadronic tensor} $W^{\mu\nu}$:
\begin{equation}
\boxed{
W^{\mu\nu}(p_B,q)
\equiv \frac{1}{4\pi m_p}\sum_X (2\pi)^4\delta^{(4)}(p_B+q-p_X)\,
\langle p|j^\mu_{\rm QCD}(0)|X\rangle
\langle p|j^\nu_{\rm QCD}(0)|X\rangle^\ast.
}
\end{equation}

\subsection{General decomposition: structure functions $W_1$ and $W_2$}
Lorentz invariance, parity invariance, and current conservation imply the standard decomposition:
\begin{equation}
\boxed{
W^{\mu\nu}(p_B,q)
= W_1(Q^2,\nu)\left(-g^{\mu\nu}+\frac{q^\mu q^\nu}{q^2}\right)
+ W_2(Q^2,\nu)\left(p_B^\mu-\frac{p_B\!\cdot q}{q^2}q^\mu\right)
\left(p_B^\nu-\frac{p_B\!\cdot q}{q^2}q^\nu\right),
}
\label{eq:hadronic-tensor-W1W2}
\end{equation}
where the energy transfer variable is
\begin{equation}
\nu \equiv \frac{p_B\cdot q}{m_p} = q^0_{\rm LAB} = E_A-E_1 \quad \text{(LAB frame)}.
\end{equation}
The scalar functions $W_1$ and $W_2$ are the \textbf{structure functions}. They play
the role that $F_1,F_2$ played in elastic scattering, but now they depend on \emph{two}
Lorentz invariants $(Q^2,\nu)$ because the final hadronic state is not fixed.

\subsection{DIS cross section in terms of $W_1$ and $W_2$}
The cross section can be written as a contraction of the leptonic tensor $L^{\mu\nu}_e$
(with known QED form) and $W_{\mu\nu}$:
\begin{equation}
\left(\frac{d\sigma}{dE_1\,d\Omega}\right)_{\rm LAB}
=
\frac{4\alpha^2 E_1^2}{q^4}
\left(
W_2(Q^2,\nu)\cos^2\frac{\theta}{2} + 2W_1(Q^2,\nu)\sin^2\frac{\theta}{2}
\right),
\end{equation}
where $\theta$ is the electron scattering angle in the LAB and $q^2\simeq -4E_AE_1\sin^2(\theta/2)$.

\subsection{Bjorken scaling: a surprise from experiment}
For scattering off an \emph{elementary} spin-$\tfrac12$ target, kinematics plus the
delta-function constraints imply that the analogs of $\nu W_2$ and $2m W_1$ depend
on a \emph{single} variable rather than on $(Q^2,\nu)$ separately. This motivates introducing
the dimensionless scaling variable
\begin{equation}
x \equiv \frac{Q^2}{2m_p \nu},
\qquad\text{equivalently}\qquad
\omega\equiv \frac{2m_p\nu}{Q^2}=\frac{1}{x}.
\end{equation}
SLAC experiments found that, at large enough $Q^2$, the proton behaves as if composed of
\textbf{quasi-free point-like constituents}: the structure functions approximately scale
(i.e.\ depend mostly on $x$) rather than on $Q^2$ and $\nu$ independently. This is the
phenomenological basis for the \textbf{parton model}.

\section{Parton model: PDFs and the Callan--Gross relation}
\subsection{Partons as quasi-free constituents}
Model the fast proton as a collection of partons $i$ (quarks, antiquarks, and possibly
other neutral constituents) carrying a fraction $x$ of the proton momentum:
\begin{equation}
p_i^\mu \simeq x\,p_B^\mu,\qquad x\in[0,1].
\end{equation}
Let $f_i(x)$ be the probability density (parton distribution function, PDF) that parton $i$
carries momentum fraction $x$.

For a spin-$\tfrac12$ parton, the parton-level kinematics implies contributions to the
proton structure functions that localize at the kinematically allowed $x$.
Summing incoherently over partons weighted by their squared electric charges $Q_i^2$ gives
\begin{equation}
F_2(x)\equiv \sum_i Q_i^2\,x\,f_i(x),
\qquad
F_1(x)\equiv \sum_i Q_i^2\,\frac12\,f_i(x),
\end{equation}
and therefore the \textbf{Callan--Gross relation}
\begin{equation}
\boxed{2x\,F_1(x)=F_2(x),}
\end{equation}
which is characteristic of spin-$\tfrac12$ constituents.
(If partons had spin $0$, one would find $F_1(x)=0$.)

In the high-energy limit ($s\gg m_p^2$) the double-differential DIS cross section can be
presented in terms of $x$ and the inelasticity variable $y$:
\begin{equation}
y \equiv \frac{p_B\cdot q}{p_B\cdot p_A},
\qquad
\frac{d\sigma}{dx\,dy}
\simeq \frac{2\pi\alpha^2}{s\,q^4}\,F_2(x)\,\bigl(1+(1-y)^2\bigr),
\end{equation}
emphasizing that DIS essentially measures $F_2(x)$, and hence particular combinations
of PDFs.

\subsection{Flavor decomposition and valence vs sea}
Assume the charged partons are quarks/antiquarks of flavors $u,d,s,\dots$. Then
\begin{equation}
\frac{F_2(x)}{x}
=\left(\frac{2}{3}\right)^2\bigl(u(x)+\bar u(x)\bigr)
+\left(\frac{1}{3}\right)^2\bigl(d(x)+\bar d(x)\bigr)
+\left(\frac{1}{3}\right)^2\bigl(s(x)+\bar s(x)\bigr)+\cdots .
\end{equation}
It is useful to separate
\begin{equation}
u(x)=u_v(x)+u_s(x),\qquad \bar u(x)=u_s(x),
\end{equation}
and similarly for $d,s$. Here:
\begin{itemize}
  \item \textbf{valence PDFs} $q_v(x)$ encode the net quantum numbers of the hadron,
  \item \textbf{sea PDFs} $q_s(x)$ arise from $q\bar q$ pair production.
\end{itemize}
At small $x$ (large energy transfer), the sea dominates; at large $x$ the valence dominates.

\subsection{Sum rules: flavor number and momentum}
PDFs satisfy powerful normalization constraints:
\begin{itemize}
  \item \textbf{Flavor number (valence) sum rules:}
  \begin{equation}
  \int_0^1 dx\,[u(x)-\bar u(x)] = \int_0^1 dx\,u_v(x) = 2,\qquad
  \int_0^1 dx\,[d(x)-\bar d(x)] = \int_0^1 dx\,d_v(x) = 1,
  \end{equation}
  for the proton $(uud)$.
  \item \textbf{Momentum sum rule:} the proton momentum is carried by all partons,
  including electrically neutral ones:
  \begin{equation}
  1=\int_0^1 dx\,x\Bigl(u+\bar u+d+\bar d+s+\bar s+g+\cdots\Bigr),
  \end{equation}
  where $g(x)$ denotes the PDF of electrically neutral partons.
\end{itemize}
Experimentally, integrating measured $F_2$ moments indicates that a substantial fraction of
the proton momentum is carried by electrically neutral partons that do not couple directly to the
photon. In QCD, these are naturally identified with \textbf{gluons}.

\subsection{Scaling violations: the footprint of QCD}
Bjorken scaling is only approximate: structure functions develop a mild dependence on $Q^2$.
In QCD this arises because quarks radiate gluons (and gluons split into $q\bar q$),
so the parton picture evolves with the resolution scale. At very small $x$ the $Q^2$ dependence
becomes stronger. This is a key bridge from QED-based DIS phenomenology to QCD dynamics.

\section{From QED phenomenology to QCD: why the strong force must be non-abelian}
DIS and hadron spectroscopy suggest the underlying strong theory must include:
\begin{enumerate}
  \item spin-$\tfrac12$ quarks with \textbf{color} and \textbf{flavor},
  \item only \textbf{color singlet} physical states (confinement),
  \item approximate \textbf{chiral symmetry} for light quarks,
  \item \textbf{spontaneous chiral symmetry breaking} to the vector subgroup,
  \item \textbf{asymptotic freedom}: quarks behave as free at large momentum transfer,
  \item and \textbf{electrically neutral constituents} carrying momentum (gluons).
\end{enumerate}
This points to a local non-abelian gauge theory: \textbf{QCD} with gauge group $SU(3)_{\rm color}$.

\subsection{Local SU(3) gauge invariance and the QCD Lagrangian}
Let $q(x)$ be a color triplet:
\[
q(x)=\begin{pmatrix}q_1(x)\\ q_2(x)\\ q_3(x)\end{pmatrix},
\qquad q(x)\to g(x)\,q(x),\qquad g(x)\in SU(3).
\]
To make $\bar q i\slashed{\partial} q$ invariant under local $g(x)$, introduce a covariant derivative
\begin{equation}
D_\mu = \partial_\mu + i g_s\, G_\mu(x),\qquad G_\mu(x)=T^a G^a_\mu(x),
\end{equation}
where $T^a$ are the $SU(3)$ generators in the fundamental representation.
The gauge field transforms as
\begin{equation}
G_\mu \to g G_\mu g^\dagger - \frac{i}{g_s}\,g\,\partial_\mu g^\dagger.
\end{equation}
Define the field strength through the commutator:
\begin{equation}
i g_s\,G_{\mu\nu} \equiv [D_\mu,D_\nu],
\qquad
G_{\mu\nu}=T^a G^a_{\mu\nu},
\qquad
G^a_{\mu\nu}=\partial_\mu G^a_\nu-\partial_\nu G^a_\mu - g_s f^{abc}G^b_\mu G^c_\nu,
\end{equation}
where $[T^a,T^b]=i f^{abc}T^c$. The Yang--Mills term is
\begin{equation}
\mathcal L_{\rm YM}= -\frac12 \mathrm{tr}(G_{\mu\nu}G^{\mu\nu})
= -\frac14 G^a_{\mu\nu}G^{a\,\mu\nu}.
\end{equation}
Therefore, for one flavor,
\begin{equation}
\mathcal L_{\rm QCD} = -\frac14 G^a_{\mu\nu}G^{a\,\mu\nu}
+\bar q\bigl(i\slashed{D}-m_q\bigr)q,
\end{equation}
and for multiple flavors $q_j$ one sums over $j=u,d,s,\dots$ with their masses $m_j$.

\paragraph{Key qualitative difference from QED.}
In QED, photons do not carry electric charge, so there are no photon--photon self-interactions.
In QCD, gluons carry color charge, so the Yang--Mills term contains \textbf{three-gluon}
and \textbf{four-gluon} interactions. This non-abelian structure is crucial for
asymptotic freedom and confinement.

\section{$e^-e^+\to$ hadrons and the number of colors}
Define the ratio
\begin{equation}
R(e^-e^+\to{\rm hadrons}) \equiv
\frac{\sigma(e^-e^+\to{\rm hadrons})}{\sigma(e^-e^+\to\mu^+\mu^-)}.
\end{equation}
At sufficiently high energies, the photon couples to quark--antiquark pairs.
A common \emph{duality} assumption is that inclusive hadron production equals inclusive
parton production:
\[
\sigma(e^-e^+\to{\rm hadrons}) \approx \sigma(e^-e^+\to q\bar q + X),
\]
and in the simplest approximation (neglecting strong interactions at high energy),
\begin{equation}
R \simeq \sum_{j\,{\rm open}} N_c\,Q_j^2,
\end{equation}
where the sum runs over quark flavors kinematically accessible, $Q_j$ is the quark charge
in units of $e$, and $N_c$ is the number of colors. Comparing to data gives strong evidence
for $N_c=3$.

\section{Heavy quarks: NRQCD, HQET, and quarkonium}
For heavy quarks (charm, bottom), $m_Q \gg \Lambda_{\rm QCD}$:
\begin{itemize}
  \item the strong coupling at the scale $m_Q$ is relatively small,
  \item heavy quarks in hadrons move non-relativistically ($v\ll 1$),
  \item effective theories become powerful: NRQCD and HQET.
\end{itemize}

\subsection{NRQCD: non-relativistic QCD for heavy quarks}
A heavy quark can be described by a Schr\"odinger field $\psi$ in the color triplet,
with an expansion in $1/m_Q$:
\begin{equation}
\mathcal L_{\rm NRQCD}
= \psi^\dagger\left(
iD_0 + \frac{\vec D^{\,2}}{2m_Q} + \vec\mu\cdot \vec B + \cdots
\right)\psi,
\qquad \vec\mu \sim \frac{g_s}{m_Q}\,\vec S,
\end{equation}
where $\vec B$ is the chromomagnetic field (a matrix in color space).
The chromomagnetic term controls spin splittings.

\subsection{HQET: heavy--light hadrons and heavy quark symmetry}
For heavy--light hadrons $H=(Q\,l)$ (one heavy quark plus light degrees of freedom),
the hadron's center of mass is close to the heavy quark position, so the heavy quark
acts approximately as a static color source. At leading order,
\begin{equation}
\mathcal L_{\rm HQET} \simeq \psi^\dagger iD_0 \psi.
\end{equation}
This yields two approximate symmetries:
\begin{itemize}
  \item \textbf{Flavor symmetry:} leading dynamics independent of $m_Q$,
  \item \textbf{Spin symmetry:} heavy-quark spin decouples at leading order.
\end{itemize}
Consequently, heavy--light hadron masses admit an expansion
\begin{equation}
M_H = m_Q + \Lambda_0^{(l)} + \mathcal O\!\left(\frac{1}{m_Q}\right),
\end{equation}
and spin partners are approximately degenerate with splitting
\begin{equation}
M_{H^\ast}-M_H = \frac{\Lambda_2^{(l)}}{m_Q} + \mathcal O\!\left(\frac{1}{m_Q^2}\right).
\end{equation}
This explains why splittings decrease for heavier quarks (bottom vs charm).

\subsection{Heavy quarkonium: bound states $(Q\bar Q)$ and gluonic decays}
For quarkonium $H=(Q\bar Q)$, both constituents are heavy and move non-relativistically about the
center of mass. A flavor-independent potential picture is often reasonable, analogous to positronium,
but with gluons instead of photons.

Decays to \emph{light hadrons} proceed via annihilation into gluons. Since QCD is invariant under
parity and charge conjugation, the number of gluons in the dominant decay depends on the state's
$J^{PC}$ and the gluon quantum numbers (gluons have $J^{PC}=1^{--}$ for the field).
A key selection rule is:
\begin{itemize}
  \item states with $J^{P+}$ can decay to \textbf{two} gluons,
  \item states with $J^{P-}$ typically require \textbf{three} gluons.
\end{itemize}
Because each extra gluon brings a factor of $\alpha_s(m_Q)$, three-gluon decays are parametrically
suppressed relative to two-gluon decays:
\begin{equation}
\frac{\Gamma\bigl((Q\bar Q)[J^{P-}]\to ggg\bigr)}
{\Gamma\bigl((Q\bar Q)[J^{P+}]\to gg\bigr)}
\sim \alpha_s(m_Q)\ll 1.
\end{equation}
For the ground states, $^1S_0$ has $J^{PC}=0^{-+}$ and $^3S_1$ has $J^{PC}=1^{--}$, so the
vector state requires three gluons while the pseudoscalar can go to two gluons, leading to strong
hierarchies of hadronic widths (up to additional numerical factors).

\section{Summary and conceptual map}
\begin{itemize}
  \item \textbf{Exclusive processes} (fixed hadronic final state) are described by \textbf{form factors},
  obtained by the most general Lorentz-covariant decomposition of current matrix elements
  plus symmetry constraints.
  \item \textbf{Elastic scattering on spin-$\tfrac12$ hadrons} needs two form factors:
  $F_1(q^2)$ and $F_2(q^2)$ (or equivalently $G_E$ and $G_M$).
  The existence of large $F_2(0)$ reflects intrinsic structure (anomalous magnetic moment).
  \item \textbf{Inclusive processes} like DIS require summing over hadronic final states, leading to
  the hadronic tensor $W^{\mu\nu}$ and \textbf{structure functions} $W_1,W_2$.
  \item \textbf{Bjorken scaling} and the \textbf{Callan--Gross relation} provide evidence that the proton
  contains \textbf{spin-$\tfrac12$ partons}.
  \item \textbf{PDF sum rules} plus DIS moments imply the presence of \textbf{electrically neutral partons}
  carrying momentum: naturally identified with \textbf{gluons}.
  \item These facts motivate \textbf{QCD} as a non-abelian $SU(3)$ gauge theory with gluon self-interactions,
  explaining (qualitatively) both asymptotic freedom and the need for color singlets.
  \item For \textbf{heavy quarks}, effective theories (NRQCD/HQET) yield approximate spin/flavor symmetries
  and explain observed degeneracy patterns and decay hierarchies.
\end{itemize}
